\begin{abstract}
Induction heads are one of the clearest mechanistic candidates for in-context copying, so they are a natural place to look for leakage-like behavior in language models. Our main question is whether empirically identified induction-like heads in \gptsmall push the model toward copying repeated prompt content even when the prompt instead asks for a diffuse or random answer. We score attention heads with an induction-head detector in TransformerLens, ablate the top-scoring heads by zeroing their \texttt{hook\_z} activations, and compare this intervention against matched random-head ablations on three tasks: a controlled copy task, a random-choice leakage task, and a small WikiText utility check. The strongest final run identifies \headpair as the top two induction-like heads. Ablating them reduces repeated-token bias on random-choice prompts from 0.0496 to 0.0211, but it also reduces normalized entropy from 0.860 to 0.819 and increases WikiText perplexity from 89.9 to 169.7. Sensitivity runs show that the conclusion is not stable across high-scoring head subsets: adding \headlsix raises repeated-token bias instead of lowering it. These results show that induction-like heads matter for both intended and unintended copying, but they do not behave like a single large causal source of leakage in \gptsmall. Any mitigation built around these circuits must therefore be evaluated jointly for behavior, entropy, and utility.
\end{abstract}
